\renewcommand{\tema}{Examen Diagnóstico - Fuerzas y Leyes de Newton}

\twocolumn
% Instrucciones
\section{Instrucciones}

\begin{instruccionbox}
\textbf{Tiempo límite: 15 minutos para 10 preguntas}

Este examen diagnóstico tiene como objetivo evaluar tus conocimientos previos sobre fuerzas, leyes de Newton y gravitación universal. Es importante que trabajes con rapidez y precisión, simulando las condiciones de un examen de admisión real.

\textbf{Aspectos evaluados:}
\begin{enumerate}
    \item \textbf{Conocimiento del tema:} Comprensión de conceptos fundamentales de mecánica
    \item \textbf{Velocidad de respuesta:} Capacidad para resolver problemas bajo presión de tiempo
    \item \textbf{Capacidad para eliminar distractores:} Habilidad para identificar respuestas incorrectas
\end{enumerate}

\textbf{Procedimiento:}
\begin{enumerate}
    \item Responde las 10 preguntas en un máximo de 15 minutos
    \item Si una pregunta toma más de 2 minutos, marca cualquier respuesta y continúa
    \item Calcula mentalmente, \textbf{sin calculadora}
    \item Anota tus respuestas en formato: 1-A, 2-C, 3-B...
    \item Al terminar, revisa las respuestas en la siguiente sección
\end{enumerate}

\textbf{Nota:} Usa $g = 10$ m/s$^2$ para la aceleración gravitacional en la Tierra.
\end{instruccionbox}

\section{Preguntas}

\begin{preguntabox}
\subsection{Pregunta 1}
Un objeto se mueve con velocidad constante en línea recta sobre una superficie horizontal. De acuerdo con la Primera Ley de Newton, ¿qué podemos afirmar sobre las fuerzas que actúan sobre el objeto?

\begin{enumerate}
    \item[A)] No hay fuerzas actuando sobre el objeto
    \item[B)] La suma vectorial de todas las fuerzas es cero
    \item[C)] Solo actúa la fuerza de gravedad
    \item[D)] Las fuerzas se cancelan parcialmente
\end{enumerate}
\end{preguntabox}

\begin{preguntabox}
\subsection{Pregunta 2}
¿Cuál es la aceleración que experimenta un cuerpo de 5 kg cuando se le aplica una fuerza neta de 20 N?

\begin{enumerate}
    \item[A)] 4 m/s$^2$
    \item[B)] 25 m/s$^2$
    \item[C)] 100 m/s$^2$
    \item[D)] 0.25 m/s$^2$
\end{enumerate}
\end{preguntabox}

\begin{preguntabox}
\subsection{Pregunta 3}
Si la masa de un objeto se duplica mientras la fuerza neta aplicada permanece constante, ¿qué sucede con su aceleración?

\begin{enumerate}
    \item[A)] Se duplica
    \item[B)] Se cuadruplica
    \item[C)] Se reduce a la mitad
    \item[D)] Permanece constante
\end{enumerate}
\end{preguntabox}

\begin{preguntabox}
\subsection{Pregunta 4}
Un libro descansa sobre una mesa. La mesa ejerce una fuerza hacia arriba sobre el libro. Según la Tercera Ley de Newton, ¿cuál es la reacción a esta fuerza?

\begin{enumerate}
    \item[A)] El peso del libro
    \item[B)] La fuerza que el libro ejerce hacia abajo sobre la mesa
    \item[C)] La fuerza que la Tierra ejerce sobre el libro
    \item[D)] La fuerza de fricción entre el libro y la mesa
\end{enumerate}
\end{preguntabox}

\begin{preguntabox}
\subsection{Pregunta 5}
¿Cuál es el peso de un objeto de 8 kg en la superficie terrestre? (Usa $g = 10$ m/s$^2$)

\begin{enumerate}
    \item[A)] 8 N
    \item[B)] 18 N
    \item[C)] 80 N
    \item[D)] 800 N
\end{enumerate}
\end{preguntabox}

\newpage

\begin{preguntabox}
\subsection{Pregunta 6}
Sobre un objeto actúan tres fuerzas horizontales: 10 N hacia la derecha, 5 N hacia la derecha y 8 N hacia la izquierda. ¿Cuál es la fuerza neta sobre el objeto?

\begin{enumerate}
    \item[A)] 23 N hacia la derecha
    \item[B)] 7 N hacia la derecha
    \item[C)] 3 N hacia la izquierda
    \item[D)] 7 N hacia la izquierda
\end{enumerate}
\end{preguntabox}

\begin{preguntabox}
\subsection{Pregunta 7}
Un objeto cuelga de dos cuerdas que forman un ángulo con la vertical. Para que el objeto esté en equilibrio traslacional, ¿qué condición debe cumplirse?

\begin{enumerate}
    \item[A)] La suma de las tensiones debe ser mayor que el peso
    \item[B)] Las componentes verticales de las tensiones deben sumar cero
    \item[C)] La suma vectorial de todas las fuerzas debe ser cero
    \item[D)] Las tensiones deben ser iguales entre sí
\end{enumerate}
\end{preguntabox}

\begin{preguntabox}
\subsection{Pregunta 8}
Una barra homogénea de 6 m de longitud está apoyada en un punto pivote. Para que esté en equilibrio rotacional con pesos iguales en ambos extremos, ¿dónde debe colocarse el pivote?

\begin{enumerate}
    \item[A)] A 2 m de un extremo
    \item[B)] A 3 m de un extremo
    \item[C)] A 4 m de un extremo
    \item[D)] En cualquier punto, no importa
\end{enumerate}
\end{preguntabox}

\newpage

\begin{preguntabox}
\subsection{Pregunta 9}
Un resorte se estira 10 cm cuando se le aplica una fuerza de 20 N. ¿Cuál es la constante elástica del resorte?

\begin{enumerate}
    \item[A)] 2 N/m
    \item[B)] 20 N/m
    \item[C)] 200 N/m
    \item[D)] 2000 N/m
\end{enumerate}
\end{preguntabox}

\begin{preguntabox}
\subsection{Pregunta 10}
Dos masas de 3 kg y 6 kg se encuentran separadas por una distancia de 2 m. Si la distancia se reduce a la mitad, ¿cómo cambia la fuerza gravitacional entre ellas?

\begin{enumerate}
    \item[A)] Se reduce a la mitad
    \item[B)] Se duplica
    \item[C)] Se cuadruplica
    \item[D)] Permanece constante
\end{enumerate}
\end{preguntabox}

\newpage
\onecolumn
\section{Respuestas y Retroalimentación}

\begin{respuestabox}
\subsection{Pregunta 1}
\textcolor{verdecorrecto}{\textbf{Respuesta correcta: B) La suma vectorial de todas las fuerzas es cero}}

\textbf{Justificación:}\\
La Primera Ley de Newton establece que un objeto permanece en reposo o en movimiento rectilíneo uniforme (velocidad constante) cuando la suma de todas las fuerzas que actúan sobre él es cero:
$$\sum \vec{F} = 0$$

Si el objeto se mueve con velocidad constante, significa que no hay aceleración ($a = 0$), y por la Segunda Ley ($\vec{F} = m\vec{a}$), la fuerza neta debe ser cero.

\textbf{Respuestas incorrectas:}
\begin{itemize}
    \item \textbf{A) No hay fuerzas actuando sobre el objeto:} Error conceptual. Siempre hay fuerzas actuando (peso, normal, posiblemente fricción), pero están equilibradas.
    \item \textbf{C) Solo actúa la fuerza de gravedad:} Incorrecto. Si solo actuara la gravedad, el objeto aceleraría hacia abajo, no se movería con velocidad constante horizontal.
    \item \textbf{D) Las fuerzas se cancelan parcialmente:} Impreciso. No se cancelan parcialmente, sino completamente (suma vectorial = 0).
\end{itemize}
\end{respuestabox}

\begin{respuestabox}

\subsection{Pregunta 2}
\textcolor{verdecorrecto}{\textbf{Respuesta correcta: A) 4 m/s$^2$}}

\textbf{Justificación:}\\
\textbf{Fórmula utilizada:} Segunda Ley de Newton
$$F = ma$$

\textbf{Despeje para $a$:}
$$a = \frac{F}{m}$$

\textbf{Sustitución:}
$$a = \frac{20 \text{ N}}{5 \text{ kg}} = 4 \text{ m/s}^2$$

\textbf{Respuestas incorrectas:}
\begin{itemize}
    \item \textbf{B) 25 m/s$^2$:} Error aritmético, suma en lugar de división: $20 + 5 = 25$.
    \item \textbf{C) 100 m/s$^2$:} Error aritmético, multiplicación en lugar de división: $20 \times 5 = 100$.
    \item \textbf{D) 0.25 m/s$^2$:} Error aritmético, inversión de la división: $5/20 = 0.25$.
\end{itemize}
\end{respuestabox}
\newpage
\begin{respuestabox}

\subsection{Pregunta 3}
\textcolor{verdecorrecto}{\textbf{Respuesta correcta: C) Se reduce a la mitad}}

\textbf{Justificación:}\\
\textbf{Fórmula utilizada:} Segunda Ley de Newton
$$a = \frac{F}{m}$$

La aceleración es inversamente proporcional a la masa cuando la fuerza es constante.

Si la masa se duplica ($m' = 2m$):
$$a' = \frac{F}{2m} = \frac{1}{2} \cdot \frac{F}{m} = \frac{a}{2}$$

Por lo tanto, la aceleración se reduce a la mitad.

\textbf{Respuestas incorrectas:}
\begin{itemize}
    \item \textbf{A) Se duplica:} Error conceptual. Confunde relación directa con inversa.
    \item \textbf{B) Se cuadruplica:} Error conceptual. Aplicaría si hubiera una relación cuadrática, que no existe aquí.
    \item \textbf{D) Permanece constante:} Error conceptual. Ignora la dependencia de $a$ con $m$.
\end{itemize}
\end{respuestabox}

\begin{respuestabox}
\subsection{Pregunta 4}
\textcolor{verdecorrecto}{\textbf{Respuesta correcta: B) La fuerza que el libro ejerce hacia abajo sobre la mesa}}

\textbf{Justificación:}\\
La Tercera Ley de Newton establece que para cada acción hay una reacción igual en magnitud, opuesta en dirección y que actúa sobre cuerpos diferentes.

Si la mesa ejerce una fuerza $\vec{F}$ hacia arriba sobre el libro (fuerza normal), el par acción-reacción es la fuerza que el libro ejerce hacia abajo sobre la mesa con magnitud $\vec{F}$.

\textbf{Criterios para identificar pares acción-reacción:}
- Misma magnitud
- Direcciones opuestas
- Actúan sobre cuerpos diferentes
- Misma naturaleza (ambas son fuerzas de contacto)

\textbf{Respuestas incorrectas:}
\begin{itemize}
    \item \textbf{A) El peso del libro:} Error conceptual. El peso actúa sobre el libro, no sobre la mesa. Además, peso y normal no forman par acción-reacción.
    \item \textbf{C) La fuerza que la Tierra ejerce sobre el libro:} Es el peso, no el par de la normal. Actúa sobre el libro, no sobre la mesa.
    \item \textbf{D) La fuerza de fricción entre el libro y la mesa:} Irrelevante para este caso. No es el par de la fuerza normal.
\end{itemize}
\end{respuestabox}
\newpage
\begin{respuestabox}

\subsection{Pregunta 5}
\textcolor{verdecorrecto}{\textbf{Respuesta correcta: C) 80 N}}

\textbf{Justificación:}\\
\textbf{Fórmula utilizada:} Peso
$$W = mg$$

donde:
- $W$ = peso (N)
- $m$ = masa (kg)
- $g$ = aceleración gravitacional (m/s$^2$)

\textbf{Sustitución:}
$$W = (8 \text{ kg})(10 \text{ m/s}^2) = 80 \text{ N}$$

\textbf{Respuestas incorrectas:}
\begin{itemize}
    \item \textbf{A) 8 N:} Error conceptual. Confunde masa con peso, usando solo el valor numérico de la masa.
    \item \textbf{B) 18 N:} Error aritmético, suma en lugar de multiplicación: $8 + 10 = 18$.
    \item \textbf{D) 800 N:} Error de unidades o de cálculo. Multiplicación incorrecta con un cero extra.
\end{itemize}
\end{respuestabox}

\begin{respuestabox}
\subsection{Pregunta 6}
\textcolor{verdecorrecto}{\textbf{Respuesta correcta: B) 7 N hacia la derecha}}

\textbf{Justificación:}\\
Para calcular la fuerza neta, sumamos vectorialmente todas las fuerzas. Establecemos un sistema de referencia donde la derecha es positiva y la izquierda es negativa.

\textbf{Fuerzas hacia la derecha (+):} $10 \text{ N} + 5 \text{ N} = 15 \text{ N}$

\textbf{Fuerzas hacia la izquierda (-):} $-8 \text{ N}$

\textbf{Fuerza neta:}
$$F_{neta} = 15 \text{ N} - 8 \text{ N} = 7 \text{ N (hacia la derecha)}$$

El signo positivo indica dirección hacia la derecha.

\textbf{Respuestas incorrectas:}
\begin{itemize}
    \item \textbf{A) 23 N hacia la derecha:} Error conceptual. Suma todas las fuerzas sin considerar direcciones: $10 + 5 + 8 = 23$.
    \item \textbf{C) 3 N hacia la izquierda:} Error aritmético. Restó incorrectamente: $8 - 5 = 3$, ignorando la fuerza de 10 N.
    \item \textbf{D) 7 N hacia la izquierda:} Magnitud correcta pero dirección incorrecta. Invirtió el signo del resultado.
\end{itemize}
\end{respuestabox}
\newpage
\begin{respuestabox}

\subsection{Pregunta 7}
\textcolor{verdecorrecto}{\textbf{Respuesta correcta: C) La suma vectorial de todas las fuerzas debe ser cero}}

\textbf{Justificación:}\\
\textbf{Primera condición de equilibrio (equilibrio traslacional):}
$$\sum \vec{F} = 0$$

Para que un objeto esté en equilibrio traslacional, la suma vectorial de todas las fuerzas que actúan sobre él debe ser cero. Esto implica que:
$$\sum F_x = 0 \quad \text{y} \quad \sum F_y = 0$$

En el caso del objeto colgado, las componentes verticales de las tensiones deben equilibrar el peso, y las componentes horizontales deben cancelarse entre sí.

\textbf{Respuestas incorrectas:}
\begin{itemize}
    \item \textbf{A) La suma de las tensiones debe ser mayor que el peso:} Error conceptual. Las tensiones en magnitud pueden ser mayores, pero sus componentes verticales deben igualar el peso, no superarlo.
    \item \textbf{B) Las componentes verticales de las tensiones deben sumar cero:} Incompleto. Deben sumar el peso (hacia arriba), no cero. La suma total de fuerzas (incluyendo el peso) es cero.
    \item \textbf{D) Las tensiones deben ser iguales entre sí:} Solo si el arreglo es simétrico. No es una condición general.
\end{itemize}
\end{respuestabox}

\begin{respuestabox}

\subsection{Pregunta 8}
\textcolor{verdecorrecto}{\textbf{Respuesta correcta: B) A 3 m de un extremo}}

\textbf{Justificación:}\\
\textbf{Segunda condición de equilibrio (equilibrio rotacional):}
$$\sum \tau = 0$$

Para una barra homogénea con pesos iguales en ambos extremos, el equilibrio rotacional se logra cuando el pivote está en el centro de masa del sistema.

Si la barra tiene 6 m de longitud y los pesos en los extremos son iguales, las torcas se equilibran cuando las distancias al pivote son iguales:

$$d_1 = d_2$$

Por lo tanto, el pivote debe estar en el centro:
$$d = \frac{6 \text{ m}}{2} = 3 \text{ m desde cualquier extremo}$$

\textbf{Respuestas incorrectas:}
\begin{itemize}
    \item \textbf{A) A 2 m de un extremo:} Genera torcas desbalanceadas. El brazo más largo produciría mayor torca.
    \item \textbf{C) A 4 m de un extremo:} Genera torcas desbalanceadas en dirección opuesta.
    \item \textbf{D) En cualquier punto, no importa:} Error conceptual. La posición del pivote determina las torcas y el equilibrio.
\end{itemize}
\end{respuestabox}
\newpage
\begin{respuestabox}

\subsection{Pregunta 9}
\textcolor{verdecorrecto}{\textbf{Respuesta correcta: C) 200 N/m}}

\textbf{Justificación:}\\
\textbf{Fórmula utilizada:} Ley de Hooke
$$F = kx$$

donde:
- $F$ = fuerza aplicada (N)
- $k$ = constante elástica (N/m)
- $x$ = elongación (m)

\textbf{Despeje para $k$:}
$$k = \frac{F}{x}$$

\textbf{Conversión de unidades:}
$$x = 10 \text{ cm} = 0.1 \text{ m}$$

\textbf{Sustitución:}
$$k = \frac{20 \text{ N}}{0.1 \text{ m}} = 200 \text{ N/m}$$

\textbf{Respuestas incorrectas:}
\begin{itemize}
    \item \textbf{A) 2 N/m:} Error de conversión. Usó centímetros sin convertir: $20/10 = 2$.
    \item \textbf{B) 20 N/m:} Error aritmético. Dividió incorrectamente o confundió la conversión.
    \item \textbf{D) 2000 N/m:} Error de conversión. Multiplicó por 10 en lugar de dividir: usó 0.01 m en lugar de 0.1 m.
\end{itemize}
\end{respuestabox}

\begin{respuestabox}

\subsection{Pregunta 10}
\textcolor{verdecorrecto}{\textbf{Respuesta correcta: C) Se cuadruplica}}

\textbf{Justificación:}\\
\textbf{Fórmula utilizada:} Ley de Gravitación Universal de Newton
$$F = G\frac{m_1 m_2}{r^2}$$

La fuerza gravitacional es inversamente proporcional al cuadrado de la distancia.

\textbf{Situación inicial:}
$$F_i = G\frac{m_1 m_2}{r^2}$$

\textbf{Situación final} (distancia reducida a la mitad: $r' = r/2$):
$$F_f = G\frac{m_1 m_2}{(r/2)^2} = G\frac{m_1 m_2}{r^2/4} = 4 \cdot G\frac{m_1 m_2}{r^2} = 4F_i$$

Por lo tanto, la fuerza se cuadruplica.

\textbf{Respuestas incorrectas:}
\begin{itemize}
    \item \textbf{A) Se reduce a la mitad:} Error conceptual. Confunde relación inversa con la dependencia cuadrática.
    \item \textbf{B) Se duplica:} Error conceptual. Ignora el exponente cuadrado en el denominador.
    \item \textbf{D) Permanece constante:} Error conceptual grave. La fuerza gravitacional depende de la distancia.
\end{itemize}
\end{respuestabox}